\subsection{A Doléans--Dade--Yen decomposition}
A general local martingale need not be locally $M^2$. An $L^1$ bound on its
running supremum neither makes the unobserved jump at first passage square
integrable nor bounds that jump deterministically. Thus an $M^2$ localizing
sequence cannot be assumed for an arbitrary local martingale.

Let $N$ be the local-martingale component of a given special decomposition.
We first need a Doléans--Dade--Yen decomposition
\[
 N=L+B,
\]
where $L$ is a local martingale with deterministically bounded jumps after
localization, and $B$ is adapted and locally of finite variation.
The construction has four steps:
\begin{enumerate}
\item A c\`adl\`ag path has finitely many jumps above any positive threshold
on a finite interval.
\item Enumerate these jumps chronologically by stopping times, obtaining an
adapted c\`adl\`ag finite-variation process.
\item Choose an exhaustive localization with integrable variation and take
predictable projections.
\item Control the compensated martingale jumps at announced times and prove
compatibility with stopping.
\end{enumerate}
The enumeration order of a nonmonotone dense skeleton cannot be used as
chronological jump order.

For the deterministic path argument, let $x$ be real-valued and c\`adl\`ag,
$c>0$, and $T<\infty$. The set
\[
 J_{c,T}=\{t\in(0,T]:c<|x_t-x_{t-}|\}
\]
is finite. Form the finite step path retaining those jumps in time order:
\[
 P_t=\sum_{s\in J_{c,T},\ s\leq t}(x_s-x_{s-}).
\]
It starts at zero, is right-continuous with left limits, and has exactly
the retained jumps, with zero jump at every other time. Its cumulative
absolute jump sum $C$ is increasing. For $0\leq u\leq v$, its real
variation satisfies the exact extended-real identity
\[
 \operatorname{eVariationOn}(P,[u,v])
 =\operatorname{ofReal}(C(v)-C(u)).
\]
The increment $C(v)-C(u)$ is the finite sum of absolute retained jumps in
$(u,v]$. This identity holds for each of the finite step paths.
Local oscillation events on the factorial skeleton give a measurable,
pathwise description of a large left jump occurring before a given time.
The infimum of the finite jump set, clamped to the horizon if it is empty,
is therefore a stopping time. When the set is nonempty, that time attains
its earliest member.

Passing to stochastic processes requires iteration of the chronological
debut, not merely measurability at each fixed time. This specialized first
debut uses factorial-skeleton oscillations and infima of finite sets;
it does not assume a generic optional debut theorem. Restrict to the portion
after the preceding stop by deterministic path operations. A countable
union over rational times describes the successor before $T$; when testing
times beyond $T$, include the jump at $T$ explicitly. Iteration produces
stopping times whose active coordinates are strictly increasing, belong
to the exact jump set including its endpoint $T$, and cover every large
left jump. After this finite set is exhausted, the coordinates equal $\top$.

For the chronological sum put
\[
 J^{(n)}_t(\omega)=1_{\{\tau_n(\omega)\leq t\}}
 (N_{\tau_n(\omega)}-N_{\tau_n(\omega)-})(\omega),
\]
with value zero when $\tau_n=\top$, and set $J=\sum_nJ^{(n)}$.
Strong adaptation of stopped values and indicators of stopping-time events
makes each coordinate strongly adapted; the countable sum is measurable.
Pathwise the enumeration becomes inactive after all large jumps have been
listed, so this sum agrees exactly with the finite step path above.
It inherits zero start, right continuity, left limits, global and local
bounded variation, and the exact cumulative-variation identity.
It copies every jump larger than $c$ in $(0,T]$ and has zero left jump at
other times. It is constant after $T$, and its finite partial sums are
eventually equal to the whole process on each path.
This constructs the finite large-jump process. Neither right-continuous
level hitting nor debut on a predictable interval algebra is being used
as a substitute for debut of a general optional graph.
